% Source: https://artofproblemsolving.com/community/c3h2005209_2018_national_mathcounts_sprint_18
% Stable identity: aops-topic-2005209-q01
How many assignments of digits $A,B,C,D$ (with $A\ne0$) are there such that the four-digit numbers $\overline{ABCD}$ in base $6$ and $\overline{ABCD}$ in base $10$ are both divisible by $5$?

A student argues: the digits are $0$--$5$, $D$ is $0$ or $5$, $A$ is $1$--$5$, and $\overline{ABCD}_6\equiv A+B+C+D\pmod 5$; so $A+B+C$ is $5$, $10$ or $15$, with $15$, $10$ and $1$ solutions respectively, giving $2(15+10+1)=52$. Find the correct count and the error in the argument.
